\section{Conclusion}
\label{sec:conclusion}
This paper formulates hybrid fault-tolerant workflow deployment under a
deadline and a deployment-dependent recovery-delay budget. The formulation
connects nominal makespan with the largest additional delay over a specified
fault-scenario set, while accounting for replica compute, provisioned
communication, passive reserve, and a recovery-delay cost proxy.

GCN-ALNS combines feasibility-first neighborhood search with supervised
resource-graph preferences and a common scenario verifier. Learned guidance
influences search decisions; it does not certify feasibility. The independent
evaluation protocol fixes learning and validation roles, selects the model
before testing, and preserves per-instance pairing and versioned deployment
artifacts.

\pending{After the complete rerun, add two or three sentences reporting the
supported per-scale effects and their uncertainty. State the limits of those
findings rather than claiming universal GCN superiority.}

Future work includes real-trace evaluation, richer failure-time coverage,
state-dependent learned guidance, and simulator profiling with incremental
or indexed dependency evaluation. Each extension should preserve an explicit
distinction between approximate search guidance and the fault model under
which feasibility is established.
